
\begin{obereflection}
\textbf{Latihan 7.1: Konsep}
1. Mengapa Triple DES menggunakan urutan Enkripsi-Dekripsi-Enkripsi (EDE) alih-alih Enkripsi-Enkripsi-Enkripsi (EEE)?
2. Jelaskan mekanisme kerja serangan \textit{Meet-in-the-Middle} pada Double DES secara sederhana.
3. Apa dampak dari penggunaan kunci lemah (\textit{weak keys}) pada algoritma DES?
4. Mengapa bit paritas pada kunci DES tidak menambah kekuatan kunci? Berapa panjang kunci efektifnya?
5. Jelaskan mengapa permutasi awal (IP) dan inversinya tidak menambah keamanan DES sama sekali.
6. Apa fungsi tumpang tindih bit pada tabel ekspansi $E$ terhadap penyebaran pengaruh satu bit?
7. Mengapa kriptanalisis linier Matsui, meskipun secara teori memerlukan lebih sedikit operasi daripada pencarian menyeluruh, tidak pernah dipakai untuk memecahkan DES dalam praktik?
8. Mengapa batas $2^{32}$ blok lebih berbahaya bagi 3DES daripada bagi AES-128?
\end{obereflection}

\begin{obereflection}
\textbf{Latihan 7.2: Hitungan dengan Tabel Bit}
Kerjakan dengan tabel pada Tabel~\ref{tab:ip-des}, Tabel~\ref{tab:ep-des},
Tabel~\ref{tab:sbox-des}, dan Tabel~\ref{tab:pc-des}.

1. Tentukan $C_0$ dan $D_0$ untuk kunci $\code{0000000000000000}$ (semua bit nol).
   Apakah kunci ini termasuk kunci lemah? Jelaskan.
2. Tentukan keluaran $S_1$ untuk masukan $\code{001011}$ dan $\code{111010}$.
3. Pada tabel ekspansi $E$, sebutkan posisi bit masukan 32 bit yang muncul
   \emph{dua kali} pada keluaran 48 bit.
4. Pada Kunci DES, posisi berapa saja yang merupakan bit paritas? Berapa banyak?
5. Berapa banyak subkunci yang dibangkitkan DES dan berapa panjang masing-masing?
6. Tentukan $\text{IP}(\code{8000000000000000})$, yaitu permutasi awal atas blok
   yang hanya bit pertamanya bernilai 1. Nyatakan jawabannya dalam heksadesimal.
7. Hitung waktu rata-rata pemecahan DES pada laju $10^{10}$ kunci per detik.
   Nyatakan jawabannya dalam hari.
8. Berapa entri yang harus disimpan penyerang pada serangan
   \textit{meet-in-the-middle} terhadap Double DES, dan berapa operasi yang
   dibutuhkan seluruhnya?
\end{obereflection}
